\documentclass{amsart}
% For dealing with references we use the comment environment
\usepackage{verbatim}
\newenvironment{reference}{\comment}{\endcomment}
%\newenvironment{reference}{}{}
\newenvironment{slogan}{\comment}{\endcomment}
\newenvironment{history}{\comment}{\endcomment}

% For commutative diagrams we use Xy-pic
\usepackage[all]{xy}

% We use 2cell for 2-commutative diagrams.
\xyoption{2cell}
\UseAllTwocells

% We use multicol for the list of chapters between chapters
\usepackage{multicol}

% This is generally recommended for better output
\usepackage{lmodern}
\usepackage[T1]{fontenc}

% For cross-file-references

% Package for hypertext links:
\usepackage{hyperref}

% For any local file, say "hello.tex" you want to link to please

% Theorem environments.
%
\theoremstyle{plain}
\newtheorem{theorem}[subsection]{Theorem}
\newtheorem{proposition}[subsection]{Proposition}
\newtheorem{lemma}[subsection]{Lemma}

\theoremstyle{definition}
\newtheorem{definition}[subsection]{Definition}
\newtheorem{example}[subsection]{Example}
\newtheorem{exercise}[subsection]{Exercise}
\newtheorem{situation}[subsection]{Situation}

\theoremstyle{remark}
\newtheorem{remark}[subsection]{Remark}
\newtheorem{remarks}[subsection]{Remarks}

\numberwithin{equation}{subsection}

% Macros
%
\def\lim{\mathop{\mathrm{lim}}\nolimits}
\def\colim{\mathop{\mathrm{colim}}\nolimits}
\def\Spec{\mathop{\mathrm{Spec}}}
\def\Hom{\mathop{\mathrm{Hom}}\nolimits}
\def\Ext{\mathop{\mathrm{Ext}}\nolimits}
\def\SheafHom{\mathop{\mathcal{H}\!\mathit{om}}\nolimits}
\def\SheafExt{\mathop{\mathcal{E}\!\mathit{xt}}\nolimits}
\def\Sch{\mathit{Sch}}
\def\Mor{\mathop{\mathrm{Mor}}\nolimits}
\def\Ob{\mathop{\mathrm{Ob}}\nolimits}
\def\Sh{\mathop{\mathit{Sh}}\nolimits}
\def\NL{\mathop{N\!L}\nolimits}
\def\CH{\mathop{\mathrm{CH}}\nolimits}
\def\proetale{{pro\text{-}\acute{e}tale}}
\def\etale{{\acute{e}tale}}
\def\QCoh{\mathit{QCoh}}
\def\Ker{\mathop{\mathrm{Ker}}}
\def\Im{\mathop{\mathrm{Im}}}
\def\Coker{\mathop{\mathrm{Coker}}}
\def\Coim{\mathop{\mathrm{Coim}}}

% Boxtimes
%
\DeclareMathSymbol{\boxtimes}{\mathbin}{AMSa}{"02}

%
% Macros for moduli stacks/spaces
%
\def\QCohstack{\mathcal{QC}\!\mathit{oh}}
\def\Cohstack{\mathcal{C}\!\mathit{oh}}
\def\Spacesstack{\mathcal{S}\!\mathit{paces}}
\def\Quotfunctor{\mathrm{Quot}}
\def\Hilbfunctor{\mathrm{Hilb}}
\def\Curvesstack{\mathcal{C}\!\mathit{urves}}
\def\Polarizedstack{\mathcal{P}\!\mathit{olarized}}
\def\Complexesstack{\mathcal{C}\!\mathit{omplexes}}
% \Pic is the operator that assigns to X its picard group, usage \Pic(X)
% \Picardstack_{X/B} denotes the Picard stack of X over B
% \Picardfunctor_{X/B} denotes the Picard functor of X over B
\def\Pic{\mathop{\mathrm{Pic}}\nolimits}
\def\Picardstack{\mathcal{P}\!\mathit{ic}}
\def\Picardfunctor{\mathrm{Pic}}
\def\Deformationcategory{\mathcal{D}\!\mathit{ef}}

\setlength{\emergencystretch}{3em}
\begin{document}
\title[Stacks Project, Tag 0C0J]{301 --- Homotopy exact sequence for proper flat families with geometrically connected reduced fibres}
\maketitle
\noindent Copyright (C) 2005--2025 Johan de Jong. Stacks Project authors. Source: \href{https://stacks.math.columbia.edu/tag/0C0J}{Tag 0C0J}.

\noindent Editorial difficulty note (not author text). Difficulty H2: The proof constructs an adjoint to pullback on finite etale covers using Stein factorization, rather than merely applying abstract exactness. A cover with one trivial geometric fibre component must be shown to equal the pullback of its Stein factor by a degree-one argument on a connected scheme..

\noindent Editorial provenance note (not author text). History: excerpt from revision \texttt{a0444\allowbreak 6e57e\allowbreak c1fbc\allowbreak 252a8\allowbreak 71afc\allowbreak ec775\allowbreak 2fb28\allowbreak 07b14}, pione.tex, lines 3601--3681. Reference presentation was adapted; original-author context and core support blocks were retained or added. The mathematical wording is unchanged. Licensed under GNU FDL 1.2 or later; no invariant sections or cover texts. See ../licenses/COPYING. Context blocks reproduce author definitions and supporting results; they are not additional counted problems.

\begin{proposition}
\label{proposition-first-homotopy-sequence}
Let $f : X \to S$ be a flat proper morphism of finite presentation whose
geometric fibres are connected and reduced. Assume $S$ is connected and
let $\overline{s}$ be a geometric point of $S$. Then there is an exact
sequence
$$
\pi_1(X_{\overline{s}}) \to \pi_1(X) \to \pi_1(S) \to 1
$$
of fundamental groups.
\end{proposition}

\begin{proof}
Let $Y \to X$ be a finite \'etale morphism. Consider the Stein factorization
$$
\xymatrix{
Y \ar[d] \ar[r] & X \ar[d] \\
T \ar[r] & S
}
$$
of $Y \to S$. By Lemma \href{https://stacks.math.columbia.edu/tag/0BUN}{lemma stein factorization etale (Tag 0BUN)}
the morphism $T \to S$ is finite \'etale. In this way we obtain
a functor $\textit{F\'Et}_X \to \textit{F\'Et}_S$.
For any finite \'etale morphism $U \to S$ a morphism
$Y \to U \times_S X$ over $X$ is the same thing as a morphism
$Y \to U$ over $S$ and such a morphism factors uniquely through
the Stein factorization, i.e., corresponds to a unique
morphism $T \to U$
(by the construction of the Stein factorization as a relative
normalization in More on Morphisms, Lemma
\href{https://stacks.math.columbia.edu/tag/03GY}{lemma stein universally closed (Tag 03GY)}
and factorization by
Morphisms, Lemma \href{https://stacks.math.columbia.edu/tag/035I}{lemma characterize normalization (Tag 035I)}).
Thus we see that the functors
$\textit{F\'Et}_X \to \textit{F\'Et}_S$ and
$\textit{F\'Et}_S \to \textit{F\'Et}_X$ are adjoints.
Note that the Stein factorization of $U \times_S X \to S$ is
$U$, because the fibres of $U \times_S X \to U$ are geometrically connected.

\medskip\noindent
By the discussion above and
Categories, Lemma \href{https://stacks.math.columbia.edu/tag/07RB}{lemma adjoint fully faithful (Tag 07RB)}
we conclude that
$\textit{F\'Et}_S \to \textit{F\'Et}_X$
is fully faithful, i.e., $\pi_1(X) \to \pi_1(S)$ is surjective
(Lemma \href{https://stacks.math.columbia.edu/tag/0BN6}{lemma functoriality galois surjective (Tag 0BN6)}).

\medskip\noindent
It is immediate that the composition
$\textit{F\'Et}_S \to \textit{F\'Et}_X \to \textit{F\'Et}_{X_{\overline{s}}}$
sends any $U$ to a disjoint union of copies of $X_{\overline{s}}$.
Hence $\pi_1(X_{\overline{s}}) \to \pi_1(X) \to \pi_1(S)$ is trivial
by Lemma \href{https://stacks.math.columbia.edu/tag/0BS8}{lemma composition trivial (Tag 0BS8)}.

\medskip\noindent
Let $Y \to X$ be a finite \'etale morphism with $Y$ connected such that
$Y \times_X X_{\overline{s}}$ contains a connected component $Z$
isomorphic to $X_{\overline{s}}$. Consider the Stein factorization $T$
as above. Let $\overline{t} \in T_{\overline{s}}$ be the point corresponding
to the fibre $Z$. Observe that $T$ is connected (as the image of a connected
scheme) and by the surjectivity above $T \times_S X$ is connected.
Now consider the factorization
$$
\pi : Y \longrightarrow T \times_S X
$$
Let $\overline{x} \in X_{\overline{s}}$ be any closed point. Note that
$\kappa(\overline{t}) = \kappa(\overline{s}) = \kappa(\overline{x})$
is an algebraically closed field.
Then the fibre of $\pi$ over $(\overline{t}, \overline{x})$ consists
of a unique point, namely the unique point $\overline{z} \in Z$
corresponding to $\overline{x} \in X_{\overline{s}}$ via the
isomorphism $Z \to X_{\overline{s}}$. We conclude that the finite
\'etale morphism $\pi$ has degree $1$ in a neighbourhood of
$(\overline{t}, \overline{x})$. Since $T \times_S X$ is connected
it has degree $1$ everywhere and we find that $Y \cong T \times_S X$.
Thus $Y \times_X X_{\overline{s}}$ splits completely.
Combining all of the above we see that
Lemmas \href{https://stacks.math.columbia.edu/tag/0BS9}{lemma functoriality galois ses (Tag 0BS9)} and
\href{https://stacks.math.columbia.edu/tag/0BTS}{lemma functoriality galois normal (Tag 0BTS)}
both apply and the proof is complete.
\end{proof}
\end{document}
